% ==============================================================
% Chapter Master: Classic Dynamic Programming
% File: sections/dp_master.tex
% ==============================================================

% \chapter{Classic Dynamic Programming}
\label{ch:dynamic-programming}

The previous chapter developed a recursive language for \emph{evaluating} consumption streams. Given an exogenous sequence $\{c_t\}$, the value function $\V_t$ was characterized as a fixed point of an aggregation operator $T$. Two facts made this work: the generalized mean admits an exact recursive decomposition, and the operator $T$ is a contraction with modulus $\beta$.

This chapter takes the next step: the consumption stream is no longer given. The agent \emph{chooses} it, subject to a budget constraint that links today's choices to tomorrow's resources. The mathematical machinery from the previous chapter survives almost intact. We replace the operator $T$ with a new operator $T^\star$ that includes a $\max$, and the same contraction argument delivers existence, uniqueness, and a constructive algorithm. What is genuinely new is the \emph{first-order condition}---the Euler equation---which characterizes optimal choices and connects dynamic programming back to the canonical CES problems of \cref{ch:canonical-problems}.

The CES structure plays a starring role. Homotheticity collapses the state space: the value function is linear in wealth, and the per-period problem reduces to a canonical static CES problem with the same KL-divergence representation developed in \cref{ch:information}. The dynamic problem is, in this sense, a sequence of static problems glued together by a single scalar---the marginal propensity to consume.

% ==============================================================
\section{The Sequence Problem}
\label{sec:dp-sequence-problem}
% ==============================================================

We begin by formulating the agent's problem in its most direct form: a single optimization over an entire sequence of choices. This formulation makes the difficulty transparent and motivates the recursive reformulation that follows.

% --------------------------------------------------------------
\subsection{Setup}
\label{subsec:dp-setup}
% --------------------------------------------------------------

A representative agent makes consumption and savings decisions over a horizon $t = 0, 1, \ldots, T$ (possibly $T = \infty$). At each date $t$, the agent observes the state $W_t$, chooses a control $c_t$, and the state evolves according to a transition rule.

\paragraph{The State.}
The state $W_t \in \mathbb{R}_+$ summarizes everything the agent needs to know about the past. In the canonical example, $W_t$ is wealth carried into period $t$, but the framework accommodates richer states---multiple assets, persistent income shocks, habit stocks---without modification.

\paragraph{The Control.}
The control $c_t \in [0, W_t]$ is consumption. Implicitly, savings is the residual $s_t = W_t - c_t$. We will sometimes write the problem with savings as the explicit choice variable; the two formulations are equivalent.

\paragraph{The Transition.}
Given a state $W_t$ and a control $c_t$, the next-period state is
\begin{equation}
W_{t+1} = R \cdot (W_t - c_t),
\label{eq:dp-transition-deterministic}
\end{equation}
where $R$ is the gross return on savings. The deterministic case sets the stage; we extend to $R_{t+1}$ random in \cref{subsec:dp-stochastic}.

\paragraph{The Objective.}
The agent ranks consumption streams using the recursive utility developed in \cref{ch:recursive-structure}:
\begin{equation}
\V_0 = \Mbar(c_0, \V_1; 1-\beta, \rho), \quad \V_t = \Mbar(c_t, \V_{t+1}; 1-\beta, \rho),
\label{eq:dp-recursive-objective}
\end{equation}
with terminal condition $\V_T = c_T$ in the finite-horizon case. The parameter $\rho = 1 - 1/\theta$ encodes the intertemporal elasticity of substitution.

\begin{calloutnote}[Why the Recursive Form]
We deliberately write the objective recursively rather than as a discounted sum. The recursive form makes two things clear. First, it accommodates Epstein--Zin preferences and other non-time-separable specifications without notational change---only the inner aggregator differs. Second, it shows that the present optimization problem is built on top of the evaluation operator $T$ from \cref{sec:fixed-points}: we now optimize \emph{over} the streams that $T$ evaluates.
\end{calloutnote}

% --------------------------------------------------------------
\subsection{The Sequential Formulation}
\label{subsec:dp-sequential}
% --------------------------------------------------------------

The agent's problem, written as a single optimization over the entire path, is:
\begin{equation}
\max_{\{c_t\}_{t=0}^{T}} \V_0 \quad \text{subject to} \quad W_{t+1} = R(W_t - c_t), \quad c_t \in [0, W_t], \quad W_0 \text{ given}.
\label{eq:sequence-problem}
\end{equation}

\paragraph{The Choice Object.}
The agent chooses an entire sequence $\{c_0, c_1, \ldots, c_T\}$ at date $0$. In the infinite-horizon case, this is an infinite-dimensional optimization. The Lagrangian has $T+1$ multipliers (one per budget constraint), and the first-order conditions are a system of $T+1$ Euler equations.

\paragraph{Lagrangian Approach.}
Forming the Lagrangian and taking first-order conditions yields, for each interior $t$:
\begin{equation}
\frac{\partial \V_0}{\partial c_t} = \mu_t,
\label{eq:lagrangian-foc}
\end{equation}
where $\mu_t$ is the multiplier on the period-$t$ budget constraint. Combining adjacent FOCs eliminates the multipliers and produces an Euler equation linking $c_t$ and $c_{t+1}$. This works, but it requires manipulating an infinite system. We seek a more economical formulation.

% --------------------------------------------------------------
\subsection{The Curse of the Sequence}
\label{subsec:dp-curse-sequence}
% --------------------------------------------------------------

The sequential formulation has three drawbacks that the recursive formulation will resolve.

\paragraph{Dimensionality.}
The choice space is infinite-dimensional. Numerical optimization over $\{c_t\}_{t=0}^{T}$ scales poorly as $T$ grows, and the limit $T \to \infty$ is not well-posed for direct numerical attack.

\paragraph{Time Inconsistency Risk.}
If preferences are not time-separable, optimizing over the full path at date $0$ may produce plans the agent will not want to execute at later dates. The recursive formulation builds in time consistency by construction: the period-$t$ problem takes $W_t$ as given and re-optimizes from there.

\paragraph{Lack of Structure.}
The sequential formulation does not exploit the recursive structure of the objective. The aggregator $\Mbar$ has the same form at every date; the budget constraint has the same form at every date; only the data ($W_0$, $R$, $\beta$) appear at the beginning. A formulation that respects this stationarity will be far more economical.

\begin{calloutimportant}[The Path Forward]
The recursive structure of the previous chapter showed that evaluating an infinite stream reduces to a single equation in a single unknown: the value function $\V$. We will now show that \emph{optimizing} over the stream reduces to a single equation as well---the Bellman equation---in which the unknown is again a value function, but now defined on the state $W$ rather than on time.
\end{calloutimportant}

% --------------------------------------------------------------
\subsection{Extension to Uncertainty}
\label{subsec:dp-stochastic}
% --------------------------------------------------------------

We close this section by indicating how the setup generalizes when returns are random. The full development is interleaved with the deterministic case in subsequent sections.

\paragraph{Stochastic Returns.}
Let $\{R_{t+1}\}$ be a sequence of random gross returns drawn from a distribution that may depend on a state variable $z_t$ (an aggregate shock, an investment opportunity index, etc.). The transition becomes:
\begin{equation}
W_{t+1} = R_{t+1} \cdot (W_t - c_t),
\label{eq:dp-transition-stochastic}
\end{equation}
and the objective uses the certainty-equivalent operator $\mu_t$ developed in \cref{sec:risk-weighted-averages}:
\begin{equation}
\V_t = \Mbar\big( c_t, \mu_t(\V_{t+1}); \, 1-\beta, \, 1 - 1/\theta \big), \qquad
\mu_t(X) = \big( \E_t[X^{1-\gamma}] \big)^{1/(1-\gamma)}.
\label{eq:dp-ez-objective}
\end{equation}
The deterministic case is recovered by $\gamma = 1/\theta$ (which collapses Epstein--Zin to time-separable CRRA) and degenerate $R_{t+1}$.

\paragraph{Information.}
Choices at date $t$ are measurable with respect to information available at date $t$, denoted $\mathcal{F}_t$. The agent knows $W_t$ and $z_t$ when choosing $c_t$; the realization $R_{t+1}$ is revealed at date $t+1$.

\begin{calloutnote}[Notation]
Throughout this chapter we use $\E_t[\cdot] = \E[\cdot \mid \mathcal{F}_t]$ for the conditional expectation and $\mu_t(\cdot)$ for the certainty equivalent. When the problem is deterministic, $\mu_t(X) = X$, and all formulas collapse to their certainty counterparts.
\end{calloutnote}

% --------------------------------------------------------------
\subsection{Summary}
\label{subsec:dp-problem-summary}
% --------------------------------------------------------------

\begin{calloutimportant}[Setup of the Optimization Problem]
\begin{enumerate}[label=(\roman*), leftmargin=2em]
    \item The agent's problem has a state $W_t$ (wealth), a control $c_t$ (consumption), and a transition $W_{t+1} = R_{t+1}(W_t - c_t)$.
    
    \item The objective is the recursive utility $\V_t = \Mbar(c_t, \mu_t(\V_{t+1}); 1-\beta, 1-1/\theta)$.
    
    \item The \textbf{sequence problem} is to choose $\{c_t\}$ to maximize $\V_0$ subject to the transition. This is infinite-dimensional and does not exploit the stationarity of the problem.
    
    \item The next section reformulates the problem \textbf{recursively}, defining a value function on states rather than on time.
\end{enumerate}
\end{calloutimportant}


% ==============================================================
\section{The Bellman Equation}
\label{sec:dp-bellman}
% ==============================================================

The recursive formulation rewrites the sequence problem as a single functional equation: the Bellman equation. The value function defined on the state space replaces the choice of an entire path, and the operator $T$ from the previous chapter acquires a $\max$ to become the Bellman operator $T^\star$. The contraction argument extends without modification, delivering existence, uniqueness, and a constructive algorithm.

% --------------------------------------------------------------
\subsection{The Principle of Optimality}
\label{subsec:dp-principle}
% --------------------------------------------------------------

The recursive reformulation rests on a simple observation due to Bellman.

\paragraph{The Statement.}
Let $V(W)$ denote the maximum value attainable from any date onward when the state is $W$. By stationarity of the environment, $V$ depends only on $W$, not on the date. Then:
\begin{equation}
V(W) = \max_{c \in [0, W]} \Mbar\Big( c, \, \mu\big( V(W') \big); \; 1-\beta, \, 1-1/\theta \Big),
\label{eq:bellman-equation}
\end{equation}
subject to $W' = R'(W - c)$, where $R'$ is the random return realized between today and tomorrow.

\paragraph{The Logic.}
Equation \cref{eq:bellman-equation} says: the optimal value today equals the best combination of (i) today's consumption $c$ and (ii) the optimal value starting from tomorrow's wealth $W'$. The agent need not commit to a path---only to today's choice $c$. Tomorrow, faced with wealth $W'$, the agent solves the same problem.

\begin{proposition}[Principle of Optimality]
\label{prop:principle-optimality}
A policy $\{c_t\}$ that solves the sequence problem \cref{eq:sequence-problem} also satisfies the Bellman equation \cref{eq:bellman-equation} at every date and state visited along the path. Conversely, a stationary policy $c = g(W)$ satisfying \cref{eq:bellman-equation} solves the sequence problem.
\end{proposition}

The formal proof uses two arguments: (i) any candidate policy can be decomposed into ``what we do today'' and ``what we do from tomorrow on,'' and (ii) the optimal value of the latter is, by definition, $V(W')$. The Bellman equation is the explicit decomposition.

\begin{calloutnote}[Time Consistency]
The principle of optimality is also a statement about \emph{time consistency}: at every date and state, the agent's continuation plan is itself optimal for the continuation problem. With time-separable preferences this is automatic; for Epstein--Zin it requires that the certainty equivalent $\mu$ be defined recursively, which is precisely the formulation we adopted in \cref{sec:epstein-zin}.
\end{calloutnote}

% --------------------------------------------------------------
\subsection{Deriving the Bellman Equation}
\label{subsec:dp-deriving-bellman}
% --------------------------------------------------------------

We give a direct derivation that emphasizes the parallel with the evaluation operator $T$ of \cref{sec:fixed-points}.

\paragraph{Starting Point.}
Define $V(W)$ as the supremum of $\V_0$ over all admissible policies, given $W_0 = W$. Then for any consumption choice $c$ today, the agent achieves at least
\begin{equation}
\Mbar\Big( c, \, \mu\big( V(R'(W - c)) \big); \; 1-\beta, \, 1-1/\theta \Big),
\label{eq:value-bound}
\end{equation}
because tomorrow's wealth is $W' = R'(W - c)$ and the agent can attain $V(W')$ from tomorrow on.

\paragraph{Optimization Over Today.}
Maximizing over $c$:
\begin{equation}
V(W) \geq \max_{c \in [0, W]} \Mbar\Big( c, \, \mu\big( V(R'(W - c)) \big); \; 1-\beta, \, 1-1/\theta \Big).
\label{eq:bellman-ge}
\end{equation}
A symmetric argument, using that any policy starting at $W$ must choose \emph{some} $c$ today, gives the reverse inequality. Equality is the Bellman equation.

\paragraph{The Evaluation Parallel.}
Compare with the evaluation operator from \cref{eq:operator-T}:
\begin{equation}
(T\V)_t = \Mbar(c_t, \V_{t+1}; 1-\beta, \rho).
\end{equation}
The Bellman equation has the same structure, with two changes:
\begin{itemize}[leftmargin=2em]
    \item The argument of $\Mbar$ on the right is $V$ evaluated at next-period wealth $W'$, not $\V$ indexed by next date.
    \item There is a $\max$ over $c$.
\end{itemize}
Both changes reflect the shift from evaluating a given stream to choosing one.

% --------------------------------------------------------------
\subsection{The Bellman Operator}
\label{subsec:dp-bellman-operator}
% --------------------------------------------------------------

We can package the Bellman equation as a fixed-point problem for an operator on a function space.

\paragraph{Definition.}
Let $\mathcal{C}$ denote the space of bounded continuous functions $V: \mathbb{R}_+ \to \mathbb{R}$ with the sup norm $\|V\| = \sup_W |V(W)|$. Define the \emph{Bellman operator} $T^\star: \mathcal{C} \to \mathcal{C}$ by
\begin{equation}
(T^\star V)(W) = \max_{c \in [0, W]} \Mbar\Big( c, \, \mu\big( V(R'(W-c)) \big); \; 1-\beta, \, 1-1/\theta \Big).
\label{eq:bellman-operator}
\end{equation}
The Bellman equation is the fixed-point condition $V = T^\star V$.

\paragraph{From $T$ to $T^\star$.}
The operator $T$ of \cref{sec:fixed-points} took a sequence-indexed value $\V_t$ and returned the same. The operator $T^\star$ takes a function $V(W)$ on the state space and returns the same. The shift in domain---from time index to state---is the substantive content of recursive optimization. The $\max$ is what distinguishes optimization from evaluation; it is the new ingredient.

\begin{calloutnote}[The State Space and Time]
In the evaluation problem, the consumption stream is given as a function of time, so the value function inherits a time index: $\V_t$. In the optimization problem, choices depend on the current state, so the value function depends on the state: $V(W)$. Time disappears as an explicit argument because the environment is stationary---the same $\Mbar$, the same $\mu$, the same $R'$ distribution at every date. This is the dimensionality reduction promised at the end of \cref{subsec:dp-curse-sequence}.
\end{calloutnote}

% --------------------------------------------------------------
\subsection{Existence, Uniqueness, and Convergence}
\label{subsec:dp-existence}
% --------------------------------------------------------------

The contraction argument from \cref{sec:fixed-points} extends almost verbatim. The presence of the $\max$ is innocuous: the maximum of two contractions is itself a contraction.

\begin{proposition}[Contraction Property of $T^\star$]
\label{prop:bellman-contraction}
The Bellman operator $T^\star$ defined in \cref{eq:bellman-operator} is a contraction with modulus $\beta$ on $(\mathcal{C}, \|\cdot\|)$. That is, for any $V_1, V_2 \in \mathcal{C}$:
\begin{equation}
\|T^\star V_1 - T^\star V_2\| \leq \beta \, \|V_1 - V_2\|.
\label{eq:contraction-Tstar}
\end{equation}
\end{proposition}

\begin{proof}[Proof sketch]
Fix $W$ and let $c_1\opt, c_2\opt$ be the maximizers in $T^\star V_1, T^\star V_2$ respectively. By optimality:
\begin{align}
(T^\star V_1)(W) - (T^\star V_2)(W) &\leq \Mbar(c_2\opt, \mu(V_1)) - \Mbar(c_2\opt, \mu(V_2)) \notag \\
&\leq \beta \, \|V_1 - V_2\|,
\end{align}
where the second inequality uses the Lipschitz property of $\Mbar$ in its second argument with modulus $\beta$ (the weight on the continuation value), together with the non-expansion of $\mu$. The reverse inequality follows by symmetry, and taking the supremum over $W$ delivers \cref{eq:contraction-Tstar}.
\end{proof}

\paragraph{Banach Once More.}
By the Banach fixed-point theorem (\cref{thm:banach}), $T^\star$ has a unique fixed point $V^\star \in \mathcal{C}$, and for any initial guess $V^{(0)} \in \mathcal{C}$, the iterates $V^{(n+1)} = T^\star V^{(n)}$ satisfy
\begin{equation}
\|V^{(n)} - V^\star\| \leq \beta^n \|V^{(0)} - V^\star\|.
\label{eq:bellman-convergence}
\end{equation}
Geometric convergence at rate $\beta$.

\paragraph{The Optimal Policy.}
Once $V^\star$ is known, the optimal policy is
\begin{equation}
g(W) = \arg\max_{c \in [0, W]} \Mbar\Big( c, \, \mu\big( V^\star(R'(W-c)) \big); \; 1-\beta, \, 1-1/\theta \Big).
\label{eq:optimal-policy}
\end{equation}
Under standard concavity and continuity assumptions on $\Mbar$ and the constraint set, $g$ is single-valued and continuous (the theorem of the maximum). This is the policy function the agent executes period by period.

\begin{calloutimportant}[What We Have Achieved]
The infinite-dimensional sequence problem of \cref{eq:sequence-problem} has become a fixed-point problem in a function space:
\[
V^\star = T^\star V^\star, \qquad g(W) = \arg\max \text{ at } V^\star.
\]
Once $V^\star$ is computed---by VFI, by backward induction in the finite-horizon case, or by exploiting structure in special cases---the entire stream $\{c_t\}$ is generated by iterating the policy $g$ on the state, starting from $W_0$. The agent's problem reduces to two objects: the function $V^\star(\cdot)$ and the function $g(\cdot)$.
\end{calloutimportant}

% --------------------------------------------------------------
\subsection{Finite Horizon and Backward Induction}
\label{subsec:dp-finite-horizon}
% --------------------------------------------------------------

When $T < \infty$, the Bellman recursion is solved exactly by backward induction.

\paragraph{Terminal Condition.}
At date $T$, the agent consumes everything: $V_T(W) = W$ (or $V_T(W) = u(W)$ with the standard CRRA period utility, as an alternative).

\paragraph{Backward Recursion.}
For $t = T-1, T-2, \ldots, 0$:
\begin{equation}
V_t(W) = \max_{c \in [0, W]} \Mbar\Big( c, \, \mu_t\big( V_{t+1}(R_{t+1}(W-c)) \big); \; 1-\beta, \, 1-1/\theta \Big).
\label{eq:backward-induction}
\end{equation}
Each step solves a static optimization over $c$ given the previously computed $V_{t+1}$. The full solution requires $T$ such optimizations, each over a one-dimensional choice variable.

\paragraph{Connection to the Infinite-Horizon Solution.}
As $T \to \infty$, $V_0(W) \to V^\star(W)$. This recovers the result of \cref{subsec:backward-finite}: the infinite-horizon fixed point is the limit of finite-horizon backward inductions.

\begin{calloutnote}[Practical Algorithm]
For numerical work, even infinite-horizon problems are solved by backward induction on a truncated horizon, with the choice of $T$ determined by the desired precision. Geometric convergence at rate $\beta$ means that $T = \log(\epsilon)/\log(\beta)$ iterations suffice for tolerance $\epsilon$. With $\beta = 0.96$ and $\epsilon = 10^{-6}$, this is roughly 340 iterations.
\end{calloutnote}

% --------------------------------------------------------------
\subsection{Summary}
\label{subsec:dp-bellman-summary}
% --------------------------------------------------------------

\begin{calloutimportant}[Key Results]
\begin{enumerate}[label=(\roman*), leftmargin=2em]
    \item The \textbf{Bellman equation} $V = T^\star V$ rewrites the sequence problem as a fixed-point problem on the state space.
    
    \item The \textbf{Bellman operator} $T^\star$ extends the evaluation operator $T$ from \cref{sec:fixed-points} by adding a $\max$ over today's control.
    
    \item $T^\star$ is a \textbf{contraction} with modulus $\beta$, inheriting this property from $T$. Existence, uniqueness, and geometric convergence follow from the Banach theorem.
    
    \item The \textbf{policy function} $g(W)$, the maximizer at $V^\star$, generates the optimal stream by iteration on the state.
    
    \item In the \textbf{finite-horizon} case, backward induction solves the problem exactly in $T$ steps. The infinite-horizon solution is the limit.
\end{enumerate}
\end{calloutimportant}

The Bellman equation gives us \emph{existence} and a constructive method, but it does not yet tell us much about the \emph{shape} of the optimal policy. For that we turn to the first-order and envelope conditions, which extract the Euler equation---the workhorse characterization of intertemporal optimality.


% ==============================================================
\section{Optimality Conditions}
\label{sec:dp-euler}
% ==============================================================

The Bellman equation establishes existence and structure of the value function. The Euler equation, derived by combining first-order and envelope conditions, characterizes optimal choices directly. It links marginal utility today to expected marginal utility tomorrow through the return on savings, and provides the bridge to asset pricing via the stochastic discount factor.

% --------------------------------------------------------------
\subsection{First-Order Conditions}
\label{subsec:dp-foc}
% --------------------------------------------------------------

We begin with the deterministic case for transparency, then extend to uncertainty.

\paragraph{The Deterministic Bellman.}
With certain return $R$, the Bellman equation is
\begin{equation}
V(W) = \max_{c \in [0, W]} \Mbar\Big( c, \, V(R(W-c)); \; 1-\beta, \, 1-1/\theta \Big).
\label{eq:bellman-det}
\end{equation}
Because $\Mbar$ is monotone in its second argument and $V$ is increasing, the right-hand side is differentiable in $c$ at an interior optimum.

\paragraph{Differentiating the Aggregator.}
Recall the canonical mean (used here in classical form for clarity):
\begin{equation}
\Mbar(x_1, x_2; \omega, \rho) = \big( \omega x_1^\rho + (1-\omega) x_2^\rho \big)^{1/\rho}.
\label{eq:agg-recall}
\end{equation}
The marginal contribution of $x_1$ to $\Mbar$ is, by the chain rule:
\begin{equation}
\frac{\partial \Mbar}{\partial x_1} = \omega \, x_1^{\rho-1} \, \Mbar^{1-\rho}.
\label{eq:dMbar-dx1}
\end{equation}
Symmetric expression for $\partial \Mbar / \partial x_2$, with $\omega$ replaced by $1-\omega$.

\paragraph{The First-Order Condition.}
Let $\rho = 1 - 1/\theta$ and $\omega = 1-\beta$. Differentiating \cref{eq:bellman-det} with respect to $c$ and setting to zero:
\begin{equation}
(1-\beta) c^{-1/\theta} \cdot \Mbar^{1/\theta} \;=\; \beta \, V(W')^{-1/\theta} \cdot \Mbar^{1/\theta} \cdot V'(W') \cdot R,
\label{eq:foc-raw}
\end{equation}
where $W' = R(W-c)$. Cancelling $\Mbar^{1/\theta}$:
\begin{equation}
(1-\beta) \, c^{-1/\theta} \;=\; \beta \, V(W')^{-1/\theta} \, V'(W') \, R.
\label{eq:foc-clean}
\end{equation}
This equates today's marginal utility to the discounted return-weighted marginal utility of next-period \emph{wealth}. To interpret the right-hand side, we need an expression for $V'(W')$. That is the role of the envelope theorem.

% --------------------------------------------------------------
\subsection{The Envelope Theorem}
\label{subsec:dp-envelope}
% --------------------------------------------------------------

The envelope theorem allows us to compute $V'(W)$ without re-optimizing, by exploiting that the optimum is interior.

\paragraph{Statement.}
Let $g(W)$ be the optimal consumption policy. Then
\begin{equation}
V(W) = \Mbar\big( g(W), \, V(R(W - g(W))); \; 1-\beta, \, 1-1/\theta \big).
\label{eq:V-at-optimum}
\end{equation}
Differentiating both sides with respect to $W$ and using the FOC (so that the indirect effect through $g'(W)$ vanishes):
\begin{equation}
V'(W) = \frac{\partial \Mbar}{\partial x_2} \cdot V'(W') \cdot R = \beta \, V(W')^{-1/\theta} \cdot V'(W') \cdot R \cdot \Mbar^{1/\theta}.
\label{eq:envelope-raw}
\end{equation}

\paragraph{Simplification via the FOC.}
Compare \cref{eq:envelope-raw} with the FOC \cref{eq:foc-clean}. The right-hand sides are identical (both equal $\beta V(W')^{-1/\theta} V'(W') R$ times $\Mbar^{1/\theta}$). Therefore:
\begin{equation}
V'(W) = (1-\beta) \, c^{-1/\theta} \, \Mbar^{1/\theta}.
\label{eq:envelope}
\end{equation}
This is the envelope condition. It says: the marginal value of wealth equals the marginal utility of consumption, scaled by the marginal contribution of consumption to the recursive aggregator.

\begin{calloutnote}[Why the Envelope Theorem Works]
The envelope theorem is the statement that, at an interior optimum, the indirect effect of changing $W$ through the policy $g(W)$ vanishes because of the FOC. So $V'(W)$ depends only on the \emph{direct} effect of $W$ on the period objective. With CES preferences, this direct effect has a particularly clean form, given by \cref{eq:envelope}.
\end{calloutnote}

% --------------------------------------------------------------
\subsection{The Euler Equation}
\label{subsec:dp-euler-derivation}
% --------------------------------------------------------------

Substituting the envelope condition \cref{eq:envelope} into the FOC \cref{eq:foc-clean} eliminates $V'$ and produces a relation purely in terms of consumption and wealth.

\paragraph{Derivation.}
At date $t$, the envelope condition gives $V'(W_t) = (1-\beta) c_t^{-1/\theta} \Mbar_t^{1/\theta}$. At date $t+1$, similarly: $V'(W_{t+1}) = (1-\beta) c_{t+1}^{-1/\theta} \Mbar_{t+1}^{1/\theta}$. Substituting the latter into the FOC:
\begin{equation}
(1-\beta) c_t^{-1/\theta} = \beta \, V(W_{t+1})^{-1/\theta} \cdot (1-\beta) c_{t+1}^{-1/\theta} \, \Mbar_{t+1}^{1/\theta} \cdot R.
\label{eq:euler-raw}
\end{equation}

Using $V(W_{t+1}) = \Mbar_{t+1}$ at the optimum (this is just the Bellman equation evaluated along the path):
\begin{equation}
c_t^{-1/\theta} = \beta R \, c_{t+1}^{-1/\theta}.
\label{eq:euler-deterministic}
\end{equation}
This is the \emph{Euler equation} in its time-separable form. The CES aggregator $\Mbar_{t+1}$ cancels out: the recursive structure leaves no residue, and we recover the familiar Euler equation of the discounted utility framework.

\paragraph{Reading the Euler Equation.}
\begin{equation}
\frac{c_{t+1}}{c_t} = (\beta R)^{\theta}.
\label{eq:euler-growth}
\end{equation}
Consumption grows when $\beta R > 1$ (return exceeds the rate of time preference) and at a rate determined by the IES $\theta$. A high IES means consumption is highly responsive to interest rates; a low IES means the agent prefers a smooth path regardless of the return.

\begin{calloutimportant}[The CES Cancellation]
The Euler equation \cref{eq:euler-deterministic} is striking in what it does \emph{not} contain: there is no trace of the recursive aggregator $\Mbar$, the value function $V$, or the path of wealth. The CES structure ensures that all of these cancel from the optimality condition, leaving a relation that involves only consumption and the return. This is the analytical advantage of CES preferences---the inter-period optimality condition is \emph{free of state variables} other than consumption and the return.
\end{calloutimportant}

% --------------------------------------------------------------
\subsection{The Stochastic Euler Equation}
\label{subsec:dp-stochastic-euler}
% --------------------------------------------------------------

When returns are random, the FOC and envelope conditions extend to expected versions, and the Euler equation acquires an expectation operator. The CES cancellation continues to hold under time-separable preferences but breaks down for Epstein--Zin, where a residual term involving the certainty-equivalent ratio survives.

\paragraph{Time-Separable Case.}
For CRRA preferences with $\gamma = 1/\theta$, the same derivation yields:
\begin{equation}
\boxed{c_t^{-1/\theta} = \beta \, \E_t\!\left[ R_{t+1} \, c_{t+1}^{-1/\theta} \right].}
\label{eq:euler-crra-stochastic}
\end{equation}
This is the canonical consumption-based asset pricing equation. In the form
\begin{equation}
1 = \E_t\!\left[ \beta \left(\frac{c_{t+1}}{c_t}\right)^{-1/\theta} R_{t+1} \right],
\label{eq:euler-asset-pricing}
\end{equation}
it says that the expected return weighted by the marginal-rate-of-substitution between today and tomorrow equals one---the absence-of-arbitrage condition.

\paragraph{The Stochastic Discount Factor.}
Define
\begin{equation}
M_{t+1} = \beta \left( \frac{c_{t+1}}{c_t} \right)^{-1/\theta}.
\label{eq:sdf-crra}
\end{equation}
Then the Euler equation becomes $\E_t[M_{t+1} R_{t+1}] = 1$, the standard pricing equation. The SDF $M_{t+1}$ is the random ``discount factor'' that prices any asset payoff. This connects dynamic programming to asset pricing: the very same first-order condition that characterizes the consumer's optimum also delivers the pricing kernel for financial markets.

\paragraph{Epstein--Zin Case.}
For general Epstein--Zin preferences ($\gamma \neq 1/\theta$), an additional term appears that captures the difference between risk aversion ($\gamma$) and intertemporal substitution ($\theta$):
\begin{equation}
\boxed{1 = \E_t\!\left[ \beta^\theta \left(\frac{c_{t+1}}{c_t}\right)^{-\theta/\theta'} \!\!\left(\frac{R_{t+1}^p \, v_{t+1}}{\mu_t(R_{t+1}^p \, v_{t+1})}\right)^{\!\theta - 1} R_{t+1} \right],}
\label{eq:euler-ez}
\end{equation}
where $\theta' = 1/(1 - 1/\gamma)$, $R_{t+1}^p$ is the return on the optimal portfolio, $v_{t+1}$ is the normalized value (developed in \cref{sec:dp-homothetic}), and $\mu_t$ is the certainty equivalent. The second factor is the Epstein--Zin correction; it vanishes when $\gamma = 1/\theta$.

\begin{calloutnote}[Two Roles of CES]
The CES structure plays two distinct roles in this section. First, in the time-separable case it makes the Euler equation \emph{simple}---the aggregator drops out. Second, in the Epstein--Zin case it makes the \emph{form} of the deviation tractable: the residual term in \cref{eq:euler-ez} is a power function of a CES ratio, not a complicated nonlinear expression. Both pieces of tractability stem from the same source: CES is the only aggregator whose marginal product is a power function.
\end{calloutnote}

% --------------------------------------------------------------
\subsection{Connection to the Static Canonical Problem}
\label{subsec:dp-static-connection}
% --------------------------------------------------------------

The Euler equation has an instructive reinterpretation as the FOC of a \emph{static} CES allocation problem.

\paragraph{The Two-Good Analogy.}
Consider a static problem: allocate one unit of resources between two goods with prices $p_1, p_2$ and CES weights $\omega, 1-\omega$:
\begin{equation}
\max_{x_1, x_2} \Mbar(x_1, x_2; \omega, \rho) \quad \text{s.t.} \quad p_1 x_1 + p_2 x_2 = 1.
\label{eq:static-ces}
\end{equation}
The FOC equates the price ratio to the marginal-rate-of-substitution:
\begin{equation}
\frac{p_1}{p_2} = \frac{\omega}{1-\omega} \left( \frac{x_1}{x_2} \right)^{\rho-1}.
\label{eq:static-foc}
\end{equation}

\paragraph{Mapping the Dynamic FOC.}
The Euler equation \cref{eq:euler-deterministic} is exactly \cref{eq:static-foc} with the substitutions:
\[
x_1 \leftrightarrow c_t, \quad x_2 \leftrightarrow c_{t+1}, \quad \omega \leftrightarrow 1-\beta, \quad \frac{p_1}{p_2} \leftrightarrow R.
\]
Today's consumption and tomorrow's consumption are two ``goods''; their relative price is the gross return $R$; the CES weights are $1-\beta$ and $\beta$. The intertemporal allocation problem is, mathematically, a two-good static CES problem.

\begin{calloutimportant}[Why This Matters]
The mapping in \cref{subsec:dp-static-connection} is not just a rhetorical flourish---it is the entry point to all of the comparative statics developed in \cref{ch:comparative-statics}. The price index, the KL representation, the substitution-versus-wealth decomposition, the elasticities of demand: \emph{all of them apply to dynamic problems with the relabeling above}. The static CES technology developed in Part I of these notes is, in this sense, a complete toolkit for the dynamic problem we are now solving.
\end{calloutimportant}

% --------------------------------------------------------------
\subsection{Summary}
\label{subsec:dp-euler-summary}
% --------------------------------------------------------------

\begin{calloutimportant}[Key Results]
\begin{enumerate}[label=(\roman*), leftmargin=2em]
    \item The \textbf{first-order condition} from the Bellman equation equates the marginal utility of consumption to the discounted return-weighted marginal value of wealth.
    
    \item The \textbf{envelope theorem} eliminates the unknown $V'$ in favor of the marginal utility of consumption: $V'(W) = (1-\beta) c^{-1/\theta} \Mbar^{1/\theta}$.
    
    \item Combining the two yields the \textbf{Euler equation} $c_t^{-1/\theta} = \beta R \, c_{t+1}^{-1/\theta}$ (deterministic) or $c_t^{-1/\theta} = \beta \E_t[R_{t+1} c_{t+1}^{-1/\theta}]$ (stochastic, time-separable).
    
    \item Defining $M_{t+1} = \beta (c_{t+1}/c_t)^{-1/\theta}$ as the \textbf{stochastic discount factor}, the Euler equation becomes the asset pricing condition $\E_t[M_{t+1} R_{t+1}] = 1$.
    
    \item The Euler equation has the same form as the FOC of a \textbf{static two-good CES} allocation problem, with consumption today and tomorrow as the two goods and the gross return as the relative price.
\end{enumerate}
\end{calloutimportant}

The Euler equation tells us how consumption should evolve along the optimal path. To convert this differential information into a level---to find $c_t$ as a function of the state $W_t$---we need to combine the Euler equation with the budget constraint. This is most cleanly done by exploiting the homothetic structure of CES preferences, which we develop next.


% ==============================================================
\section{Homothetic Structure and Wealth Normalization}
\label{sec:dp-homothetic}
% ==============================================================

CES preferences are homothetic: scaling all consumption levels by a constant scales utility by the same constant. In the dynamic problem, this property has a powerful consequence---the value function is linear in wealth, the optimal policy is a constant fraction of wealth, and the state space collapses. The per-period problem becomes a static CES problem of exactly the kind studied in \cref{ch:canonical-problems}, complete with its KL representation. This section develops the implication chain.

% --------------------------------------------------------------
\subsection{Homotheticity of CES Preferences}
\label{subsec:dp-homotheticity}
% --------------------------------------------------------------

Recall that the recursive utility $\V_t = \Mbar(c_t, \V_{t+1}; 1-\beta, \rho)$ inherits the homogeneity of $\Mbar$.

\paragraph{Homogeneity of the Aggregator.}
The classical mean $\Mbar(x_1, x_2; \omega, \rho) = (\omega x_1^\rho + (1-\omega) x_2^\rho)^{1/\rho}$ is homogeneous of degree one:
\begin{equation}
\Mbar(\lambda x_1, \lambda x_2; \omega, \rho) = \lambda \, \Mbar(x_1, x_2; \omega, \rho), \quad \lambda > 0.
\label{eq:mbar-homogeneous}
\end{equation}
Iterating, the recursive utility $\V_0$ is homogeneous of degree one in the entire consumption stream: scaling $\{c_t\}$ by $\lambda$ scales $\V_0$ by $\lambda$.

\paragraph{Homogeneity of the Constraint Set.}
The transition $W' = R(W - c)$ is also linear, so the feasible set of consumption streams scales linearly with $W_0$: from initial wealth $\lambda W_0$, the agent can attain exactly $\lambda$ times any stream attainable from $W_0$. Combined with homogeneity of preferences, this delivers the following.

\begin{proposition}[Homogeneity of the Value Function]
\label{prop:V-homogeneous}
The optimal value function $V$ in the deterministic problem (and in the stochastic problem with i.i.d.\ returns or, more generally, returns whose distribution is independent of $W$) satisfies
\begin{equation}
V(\lambda W) = \lambda V(W), \quad \lambda > 0.
\label{eq:V-homogeneous}
\end{equation}
\end{proposition}

\begin{proof}[Proof sketch]
By induction on the finite-horizon problem and passage to the limit, or directly via the principle of optimality: if $\{c_t^\star\}$ is optimal from $W$, then $\{\lambda c_t^\star\}$ is feasible from $\lambda W$, and homogeneity of preferences gives a value of $\lambda V(W)$. The optimality of $\{\lambda c_t^\star\}$ from $\lambda W$ follows because any improvement could be scaled down to improve $\{c_t^\star\}$, contradicting its optimality.
\end{proof}

% --------------------------------------------------------------
\subsection{The Normalized Bellman Equation}
\label{subsec:dp-normalized-bellman}
% --------------------------------------------------------------

Homogeneity allows us to write $V(W) = v \cdot W$ for some scalar $v$, eliminating wealth from the Bellman equation.

\paragraph{Substitution.}
Using $V(W) = v W$ and $V(W') = v W'$ in the Bellman equation \cref{eq:bellman-equation}:
\begin{equation}
v W = \max_{c \in [0, W]} \Mbar\Big( c, \, \mu(v R'(W - c)); \; 1-\beta, \, 1-1/\theta \Big).
\label{eq:bellman-vW}
\end{equation}

Pull out the random gross return: $\mu(v R'(W-c)) = (W-c) \, v \, \mu(R')$ if $v$ and $W-c$ are deterministic and $\mu$ is positively homogeneous. (The certainty equivalent $\mu(X) = (\E[X^{1-\gamma}])^{1/(1-\gamma)}$ is indeed homogeneous of degree one.) Then:
\begin{equation}
v W = \max_{c} \Mbar\Big( c, \, (W-c) \, v \, \mu(R'); \; 1-\beta, \, 1-1/\theta \Big).
\label{eq:bellman-pulled}
\end{equation}

\paragraph{Re-Parameterizing.}
Let $m = c/W$ denote the consumption-to-wealth ratio. Then $c = mW$ and $W - c = (1-m)W$. Substituting and using homogeneity of $\Mbar$:
\begin{equation}
v W = W \cdot \max_{m \in [0, 1]} \Mbar\Big( m, \, (1-m) \, v \, \mu(R'); \; 1-\beta, \, 1-1/\theta \Big).
\label{eq:bellman-m}
\end{equation}

Dividing by $W$ delivers the \emph{normalized Bellman equation}:
\begin{equation}
\boxed{v = \max_{m \in [0, 1]} \Mbar\Big( m, \, (1-m) \, v \, \mu(R'); \; 1-\beta, \, 1-1/\theta \Big).}
\label{eq:normalized-bellman}
\end{equation}
A scalar fixed-point equation in the scalar unknown $v$. Wealth has disappeared.

\begin{calloutimportant}[The State-Space Collapse]
The state variable $W$ has been eliminated. The Bellman equation, originally a functional equation for $V(\cdot)$ on $\mathbb{R}_+$, has become a single algebraic equation in a single number $v$. This is the practical payoff of CES homotheticity: a problem that would generically be infinite-dimensional is solved by characterizing one scalar.
\end{calloutimportant}

\paragraph{Optimal Consumption Rate.}
The maximizer of \cref{eq:normalized-bellman}, denoted $m^\star$, is the optimal consumption-to-wealth ratio. The first-order condition is:
\begin{equation}
(1-\beta) \, m^{-1/\theta} = \beta \, v \, \mu(R') \cdot \big( (1-m) v \mu(R') \big)^{-1/\theta}.
\label{eq:m-foc}
\end{equation}
Rearranging:
\begin{equation}
\frac{m^\star}{1 - m^\star} = \left( \frac{1-\beta}{\beta} \right)^\theta \cdot \big( v \, \mu(R') \big)^{\theta - 1}.
\label{eq:m-formula}
\end{equation}
Substituting back into \cref{eq:normalized-bellman} pins down $v$.

\paragraph{Closed Form.}
After algebra, the stationary normalized value satisfies
\begin{equation}
v = \left( (1-\beta)^\theta + \beta^\theta \mu(R')^{\theta - 1} \right)^{1/(\theta - 1)} \cdot (1-\beta)^{\theta/(1-\theta)},
\label{eq:v-closed}
\end{equation}
and the consumption rate is
\begin{equation}
m^\star = 1 - \beta^\theta \mu(R')^{\theta - 1} (1-\beta)^{1-\theta}.
\label{eq:m-closed}
\end{equation}
For $\theta = 1$ (log utility), \cref{eq:m-closed} simplifies to the classic result $m^\star = 1 - \beta$: the agent consumes a constant fraction $1-\beta$ of wealth each period, regardless of the return.

% --------------------------------------------------------------
\subsection{Recovering the Static Canonical Problem}
\label{subsec:dp-static-recovery}
% --------------------------------------------------------------

The normalized Bellman equation \cref{eq:normalized-bellman} is, structurally, the static two-good CES problem of \cref{ch:canonical-problems}.

\paragraph{The Mapping.}
Compare the canonical static problem
\begin{equation}
U^\star = \max_{x_1 + x_2 = 1} \Mbar(x_1, x_2; \omega, \rho)
\label{eq:static-canonical}
\end{equation}
(after normalizing the budget to one) with the normalized Bellman \cref{eq:normalized-bellman}, which fixes the resource as $1$ (since $m + (1-m) = 1$). The mapping is:
\[
x_1 \leftrightarrow m, \quad x_2 \leftrightarrow 1 - m, \quad \omega \leftrightarrow 1 - \beta,
\]
together with a shadow price $p_1 = 1$ on consumption today and a shadow price $p_2 = 1/(v \mu(R'))$ on consumption tomorrow.

\paragraph{The Shadow Price of the Future.}
The shadow price of next-period ``goods'' is
\begin{equation}
p_2 = \frac{1}{v \, \mu(R')}.
\label{eq:shadow-price}
\end{equation}
This is the cost (in today's units) of one unit of certainty-equivalent normalized future wealth. Three pieces enter: the inverse of the certainty equivalent of returns ($1/\mu(R')$), the inverse of the per-unit value of future wealth ($1/v$), and the implicit normalization $p_1 = 1$.

\begin{calloutnote}[The Dynamic Price Index]
The shadow price $p_2$ plays the role of the price index $\Pstar$ from \cref{ch:comparative-statics}, but for the dynamic problem. It aggregates the random return and the per-unit continuation value into a single number that prices ``the future.'' The KL representation of welfare developed in \cref{sec:price-index-kl} carries over directly: the agent's value $v$ depends on the dispersion of $\mu(R')$ across states, with the dispersion penalty governed by KL divergence.
\end{calloutnote}

% --------------------------------------------------------------
\subsection{The KL Representation in Dynamic Settings}
\label{subsec:dp-kl-representation}
% --------------------------------------------------------------

The KL representation of CES welfare (\cref{ch:information}) extends to dynamic problems with no modification beyond the relabeling above.

\paragraph{Recap: The Static Result.}
For a static CES problem with prices $\bp$ and weights $\bom$, the welfare loss from price dispersion is governed by the KL divergence between the optimal expenditure-share weights $\bw^\star$ and the preference weights $\bom$:
\begin{equation}
\frac{d \log \Pstar}{d \tilde{\rho}} = \frac{\KL{\bw^\star}{\bom}}{\tilde{\rho}^2}, \qquad \tilde{\rho} = \frac{\rho}{\rho - 1}.
\label{eq:kl-static}
\end{equation}
Higher KL divergence means greater welfare sensitivity to substitutability.

\paragraph{Two-Period Dynamic Case.}
For the two-period problem (today vs.\ tomorrow), the KL divergence is between the weights on $c_t$ and $c_{t+1}$ implied by the budget at the optimum, and the preference weights $(1-\beta, \beta)$. With deterministic return $R$:
\begin{equation}
w_1^\star = \frac{c_t}{c_t + c_{t+1}/R}, \quad w_2^\star = 1 - w_1^\star, \quad \omega_1 = 1-\beta, \quad \omega_2 = \beta.
\label{eq:dynamic-kl-weights}
\end{equation}
The KL divergence $\KL{\bw^\star}{\bom}$ measures the gap between optimal allocation shares and preference shares; it vanishes when $\beta R = 1$ (Euler equation balanced) and grows otherwise.

\paragraph{Stochastic Case.}
With random returns, the KL representation includes an additional layer: the divergence between the risk-neutral weights induced by the certainty equivalent and the preference weights. The decomposition
\begin{equation}
\frac{d \log v}{d \tilde\rho} = \KL{\text{intertemporal}}{\text{preferences}} + \KL{\text{across states}}{\text{probabilities}}
\label{eq:kl-decomposition-dyn}
\end{equation}
separates the welfare cost of price dispersion across time from the welfare cost of return dispersion across states. The first term governs the IES; the second governs risk aversion. They are decoupled in Epstein--Zin precisely because the inner and outer aggregators have different parameters.

\begin{calloutimportant}[Unification]
Three observations come together at this point:
\begin{enumerate}[leftmargin=2em]
    \item The dynamic problem is solved by a normalized Bellman equation with no state variable.
    \item The normalized Bellman equation is, structurally, a static two-good CES problem.
    \item All comparative statics, KL representations, and welfare formulas from the static chapter apply with appropriate relabeling.
\end{enumerate}
The dynamic CES problem is, in a precise sense, the static CES problem repeated. This is the theoretical basis for the applications developed in subsequent chapters---consumption and savings, portfolio choice, and labor supply---each of which is, internally, a static CES allocation.
\end{calloutimportant}

% --------------------------------------------------------------
\subsection{When Homotheticity Fails}
\label{subsec:dp-non-homothetic}
% --------------------------------------------------------------

Several extensions break homotheticity and require the full state-space treatment.

\paragraph{Borrowing Constraints.}
A constraint $W - c \geq \underline{a}$ for some absolute floor $\underline{a} > 0$ violates the linear scaling of the constraint set. Wealth re-enters as a state variable, the consumption rate $m$ varies with $W$, and the value function is no longer linear.

\paragraph{Subsistence Consumption.}
A specification $u(c - \bar{c})$ for some subsistence level $\bar{c} > 0$ breaks homogeneity at low wealth levels. The agent's effective discount rate and IES vary with wealth.

\paragraph{Labor Income.}
With non-zero labor income $y_t$, total resources are $W_t + y_t/R + \cdots$, and the relevant ``wealth'' is human plus financial. If labor income is bounded away from zero, the same homothetic argument applies to total wealth, but only in the long-horizon limit.

\paragraph{Time-Varying Investment Opportunities.}
If the distribution of returns depends on a state $z_t$ (e.g., a stochastic volatility process), the value function becomes $V(W, z) = v(z) \cdot W$ with $v$ a non-trivial function of $z$. Homotheticity in $W$ survives, but a separate fixed-point problem in $v(z)$ remains.

\begin{calloutnote}[Recovering Tractability]
In each of these cases, the KL representation and the canonical-problem mapping continue to hold \emph{conditional on the additional state}. The role of CES is to provide a one-dimensional reduction in $W$, leaving the additional state(s) as the only sources of complexity. This is the divide-and-conquer logic that makes CES-based dynamic models tractable even when full homotheticity fails.
\end{calloutnote}

% --------------------------------------------------------------
\subsection{Summary}
\label{subsec:dp-homothetic-summary}
% --------------------------------------------------------------

\begin{calloutimportant}[Key Results]
\begin{enumerate}[label=(\roman*), leftmargin=2em]
    \item CES preferences and a linear budget constraint imply a \textbf{homogeneous value function}: $V(W) = v \cdot W$ for a scalar $v$.
    
    \item The Bellman equation reduces to a \textbf{normalized Bellman equation} in the consumption rate $m$ and the per-unit value $v$, with $W$ eliminated.
    
    \item Closed-form solutions: $m^\star = 1 - \beta^\theta \mu(R')^{\theta-1} (1-\beta)^{1-\theta}$ and a corresponding closed form for $v$.
    
    \item The normalized problem has \textbf{the same structure as a static two-good CES problem}, with the gross return playing the role of an inverse price.
    
    \item The \textbf{KL representation} of welfare from \cref{ch:information} carries over to the dynamic problem, with separate KL terms for intertemporal and cross-state dispersion in the Epstein--Zin case.
    
    \item Homotheticity fails under borrowing constraints, subsistence consumption, or non-scalar wealth, but the canonical-problem mapping survives \emph{conditionally} on additional state variables.
\end{enumerate}
\end{calloutimportant}


% ==============================================================
\section{Solution Methods}
\label{sec:dp-solution}
% ==============================================================

The Bellman equation is a fixed-point problem; the Euler equation is a functional equation in the policy. Each suggests an algorithm: value function iteration (VFI) for the former, policy function iteration (PFI) for the latter. We develop both, discuss their relative merits, and conclude with a two-step decomposition that exploits CES structure to reduce the dimensionality of the inner problem.

% --------------------------------------------------------------
\subsection{Value Function Iteration}
\label{subsec:dp-vfi}
% --------------------------------------------------------------

VFI solves the Bellman equation $V = T^\star V$ by direct iteration on the operator $T^\star$.

\paragraph{The Algorithm.}
Initialize $V^{(0)}$ (e.g., $V^{(0)} \equiv 0$ or $V^{(0)} = u$). At iteration $n$:
\begin{equation}
V^{(n+1)}(W) = \max_{c \in [0, W]} \Mbar\Big( c, \, \mu\big( V^{(n)}(R'(W-c)) \big); \; 1-\beta, \, 1-1/\theta \Big).
\label{eq:vfi-step}
\end{equation}
Iterate until $\|V^{(n+1)} - V^{(n)}\| < \epsilon$.

\paragraph{Convergence.}
By the contraction property of $T^\star$ (\cref{prop:bellman-contraction}):
\begin{equation}
\|V^{(n)} - V^\star\| \leq \beta^n \|V^{(0)} - V^\star\|.
\label{eq:vfi-convergence}
\end{equation}
Geometric convergence at rate $\beta$. With $\beta = 0.96$, achieving $\epsilon = 10^{-6}$ requires about 340 iterations from any starting point.

\paragraph{Implementation Considerations.}
Each iteration requires solving \cref{eq:vfi-step} at every grid point. The cost decomposes as:
\begin{itemize}[leftmargin=2em]
    \item \textbf{Discretization:} Choose a grid $\{W_1, \ldots, W_N\}$ over the state space.
    \item \textbf{Interpolation:} Evaluating $V^{(n)}(R'(W-c))$ at off-grid values requires interpolation (linear, cubic spline, or Chebyshev).
    \item \textbf{Inner optimization:} For each $W_i$, solve a one-dimensional optimization over $c$ (golden-section, Brent's method, or grid search).
    \item \textbf{Expectation:} If returns are stochastic, $\mu(V(R'\cdot))$ must be approximated by quadrature or Monte Carlo.
\end{itemize}
Total cost per iteration is $O(N \cdot N_c \cdot N_R)$ where $N$ is the state grid, $N_c$ the number of candidate consumption levels, and $N_R$ the number of return draws.

\begin{calloutnote}[Why VFI Is the Default]
VFI is robust: it converges from \emph{any} starting point, requires no derivatives, and handles non-smooth or constrained problems gracefully. The downside is the cost per iteration. For smooth, interior problems, the methods below converge faster.
\end{calloutnote}

% --------------------------------------------------------------
\subsection{Policy Function Iteration}
\label{subsec:dp-pfi}
% --------------------------------------------------------------

PFI exploits the fact that, given a candidate policy, the value function it induces is computed cheaply (no max, just evaluation), and then the policy can be improved.

\paragraph{The Algorithm.}
Initialize a policy $g^{(0)}: \mathbb{R}_+ \to \mathbb{R}_+$. At iteration $n$:

\textbf{Step 1 (Policy Evaluation).} Compute the value $V^{(n)}$ induced by policy $g^{(n)}$, by solving the linear functional equation
\begin{equation}
V^{(n)}(W) = \Mbar\Big( g^{(n)}(W), \, \mu\big( V^{(n)}(R'(W - g^{(n)}(W))) \big); \; 1-\beta, \, 1-1/\theta \Big).
\label{eq:pfi-eval}
\end{equation}
This is a fixed-point problem in $V^{(n)}$ \emph{for fixed} $g^{(n)}$. Because no $\max$ is involved, the operator on the right is the evaluation operator $T$ from \cref{sec:fixed-points}, and the contraction argument applies. In practice, \cref{eq:pfi-eval} is often solved iteratively (a few sub-iterations of $T$) or, in special cases, analytically.

\textbf{Step 2 (Policy Improvement).} Given $V^{(n)}$, compute the new policy:
\begin{equation}
g^{(n+1)}(W) = \arg\max_{c \in [0, W]} \Mbar\Big( c, \, \mu\big( V^{(n)}(R'(W - c)) \big); \; 1-\beta, \, 1-1/\theta \Big).
\label{eq:pfi-improve}
\end{equation}

Iterate until $\|g^{(n+1)} - g^{(n)}\| < \epsilon$.

\paragraph{Howard Improvement.}
\Cref{eq:pfi-eval} is itself a fixed-point problem. If we solve it exactly at each $n$, the algorithm is \emph{exact policy iteration} and converges in $O(\log(\epsilon))$ outer iterations---superlinear convergence. If we instead apply $T$ a fixed finite number of times (say, $K$) before improving, we obtain \emph{Howard improvement}, also called \emph{modified policy iteration}. With $K = 1$, Howard improvement reduces to VFI; with $K = \infty$, to exact PFI. Intermediate $K$ trades inner cost for outer iterations.

\begin{calloutnote}[Why PFI Is Faster When It Works]
PFI converges quadratically near the optimum (as a Newton method on the Bellman equation), versus geometrically for VFI. The cost is the policy evaluation step, which itself requires solving a fixed-point problem. When this inner problem is cheap---e.g., when policies admit closed-form value functions, as in linear-quadratic problems or our homothetic CES setting---PFI dominates VFI in total compute time.
\end{calloutnote}

% --------------------------------------------------------------
\subsection{The Two-Step Decomposition for CES Problems}
\label{subsec:dp-two-step}
% --------------------------------------------------------------

For homothetic CES problems, the structure developed in \cref{sec:dp-homothetic} suggests an algorithm that exploits the static-canonical-problem mapping at every iteration.

\paragraph{The Decomposition.}
The Bellman equation in normalized form \cref{eq:normalized-bellman} factors into:
\begin{enumerate}[leftmargin=2em]
    \item An \emph{inner problem}: compute the certainty equivalent $\mu(R')$ of returns.
    \item An \emph{outer problem}: solve the two-good static CES problem in $m$ vs.\ $1-m$, with shadow price $1/(v \mu(R'))$ on the future.
\end{enumerate}
The outer problem is one-dimensional (in $m$); the inner problem is one-dimensional (in the random variable $R'$).

\paragraph{The Algorithm.}
Initialize $v^{(0)}$. At iteration $n$:

\textbf{Step 1 (Inner: Certainty Equivalent).} Compute
\begin{equation}
\mu^{(n)} = \mu(R'),
\label{eq:inner-CE}
\end{equation}
which is independent of $v^{(n)}$ and so is computed once (or at the optimum of an outer portfolio problem in the multi-asset case).

\textbf{Step 2 (Outer: Static CES Update).} Solve the two-good static problem
\begin{equation}
v^{(n+1)} = \max_{m \in [0,1]} \Mbar\big( m, \, (1-m) v^{(n)} \mu^{(n)}; \; 1-\beta, \, 1-1/\theta \big).
\label{eq:outer-CES}
\end{equation}
Use the closed-form static solution from \cref{ch:canonical-problems}: with shadow prices $(1, 1/(v^{(n)} \mu^{(n)}))$ and weights $(1-\beta, \beta)$, the optimal $m$ and the resulting value follow directly.

Iterate to convergence in $v$. Convergence is geometric, with modulus governed by the elasticity of $v$ to $v$ in \cref{eq:outer-CES}, typically much smaller than $\beta$.

\begin{calloutimportant}[Why This Works]
Each iteration in the two-step algorithm reduces to:
\begin{itemize}[leftmargin=2em]
    \item A static portfolio problem (Step 1) of the kind solved in \cref{ch:canonical-problems}.
    \item A static intertemporal allocation problem (Step 2) of the kind solved in the same chapter.
\end{itemize}
The full toolbox of comparative statics, KL representations, and closed-form solutions from Part I applies at every iteration. The dynamic problem is decomposed into a sequence of static problems for which we already have analytical solutions.
\end{calloutimportant}

\paragraph{Multi-Asset Extension.}
With multiple risky assets and portfolio weights $\balpha$, Step 1 becomes the static portfolio problem
\begin{equation}
\balpha^\star = \arg\max_{\balpha} \mu\big( R'(\balpha) \big), \qquad R'(\balpha) = \sum_j \alpha_j R_j',
\label{eq:portfolio-step}
\end{equation}
which is itself a CES problem (in the certainty equivalent form). Step 2 is unchanged. The two-step decomposition isolates the portfolio choice from the consumption-savings choice---a separation theorem that holds because of homotheticity and the time-separable risk-time structure of Epstein--Zin.

% --------------------------------------------------------------
\subsection{Projection and Collocation}
\label{subsec:dp-projection}
% --------------------------------------------------------------

When the homothetic reduction does not apply---e.g., with borrowing constraints---the value function depends nontrivially on $W$, and we need a function approximation method.

\paragraph{Projection Methods.}
Approximate $V$ by a finite-dimensional family:
\begin{equation}
V(W) \approx \hat V(W; \boldsymbol{\theta}) = \sum_{k=1}^{K} \theta_k \phi_k(W),
\label{eq:projection}
\end{equation}
where $\{\phi_k\}$ are basis functions (Chebyshev polynomials, splines, neural networks) and $\boldsymbol\theta = (\theta_1, \ldots, \theta_K)$ are coefficients.

\paragraph{Collocation.}
Choose $K$ collocation points $\{W_i\}_{i=1}^{K}$. Require the Bellman equation to hold exactly at these points:
\begin{equation}
\hat V(W_i; \boldsymbol\theta) = (T^\star \hat V)(W_i; \boldsymbol\theta), \quad i = 1, \ldots, K.
\label{eq:collocation}
\end{equation}
This is a system of $K$ equations in the $K$ coefficients $\boldsymbol\theta$. Solve by Newton iteration, or by combining with VFI (alternating updates of $\boldsymbol\theta$ via least squares and updates of the maximizers).

\paragraph{When to Use Projection.}
Projection methods are most attractive when the value function is smooth (Chebyshev), when the state has multiple dimensions (where grid methods scale poorly), or when high accuracy is required. Modern implementations using neural networks scale to very high-dimensional state spaces but at the cost of analytical transparency.

\begin{calloutnote}[The Curse of Dimensionality]
For low-dimensional problems ($d \leq 3$), VFI on a grid is reliable and easy to implement. For $d \geq 4$, the grid grows exponentially ($N^d$ points), and projection methods or simulation-based methods (e.g., simulated method of moments, deep reinforcement learning) become necessary. The CES homothetic reduction is one of the few generic tools that reduces dimensionality \emph{exactly}, without approximation. This is what makes it valuable.
\end{calloutnote}

% --------------------------------------------------------------
\subsection{Summary}
\label{subsec:dp-solution-summary}
% --------------------------------------------------------------

\begin{calloutimportant}[Key Results]
\begin{enumerate}[label=(\roman*), leftmargin=2em]
    \item \textbf{Value function iteration} (VFI) iterates the Bellman operator $T^\star$. Convergence is geometric at rate $\beta$. Robust but expensive.
    
    \item \textbf{Policy function iteration} (PFI) alternates policy evaluation and policy improvement. Converges quadratically near the optimum, but requires solving a sub-fixed-point in the evaluation step.
    
    \item \textbf{Howard improvement} interpolates between VFI and PFI, applying $T$ a finite number of times before improving. The $K=1$ case recovers VFI.
    
    \item For \textbf{homothetic CES problems}, the two-step decomposition reduces each iteration to a static portfolio problem (inner) and a static intertemporal allocation (outer). Both have closed-form solutions from \cref{ch:canonical-problems}.
    
    \item When homotheticity fails, \textbf{projection methods} approximate $V$ by a basis expansion and impose the Bellman equation at collocation points. This handles multi-dimensional states and smooth problems efficiently.
\end{enumerate}
\end{calloutimportant}

This concludes our development of the foundational tools of dynamic programming. The chapters that follow---consumption and savings, portfolio choice, labor and leisure, investment with adjustment costs---each apply this machinery to a specific economic environment. In each case, the same template recurs: write the Bellman equation, exploit homotheticity to normalize, recognize the per-period problem as a static CES allocation, and solve. The tools developed here are the universal solvent.

